\documentclass[10pt,a4paper]{amsart}
\usepackage[margin=19mm]{geometry}
\usepackage{amsmath,amssymb,mathtools}
\usepackage[scr=rsfs,cal=euler]{mathalfa}
\usepackage{xfrac}
\usepackage[unicode,hidelinks,bookmarks=false]{hyperref}
\newcommand{\ie}{i.e.\ }
\newcommand{\cf}{cf.\ }
\newcommand{\resp}{respectively\ }
\newcommand{\Subsection}{\subsection}
\providecommand{\nobreakdash}{\nobreak\hbox{-}}
\setlength{\emergencystretch}{2em}
\multlinegap0pt
\usepackage{amssymb}
\usepackage{booktabs}
\usepackage[shortlabels]{enumitem}
\usepackage{tikz}
\newtheorem{theorem}{Theorem}
\title{Electromagnetic Modes in Spherical Cavities: Complete Theory of Angular Spectra, Dispersion Relations, and Self-Adjoint Extensions}
\author{Mustafa Bakr, Tongyu Zhang,, Smain Amari}
\date{Source: arXiv:2512.18498v1 (CC BY 4.0)}
\begin{document}
\maketitle

\noindent\emph{Source and licence.} Electromagnetic Modes in Spherical Cavities: Complete Theory of Angular Spectra, Dispersion Relations, and Self-Adjoint Extensions. Version arXiv:2512.18498v1 (CC BY 4.0), distributed by arXiv under the Creative Commons Attribution 4.0 International licence, http://creativecommons.org/licenses/by/4.0/, as recorded in the arXiv licence metadata for this exact version and shown on the abstract page https://arxiv.org/abs/2512.18498v1. Full licence notice: \texttt{licenses/arxiv-2512.18498-CC-BY-4.0.txt}.\par\smallskip

\noindent\textit{Editorial scope.} Counted unit: the article's theorem thm:sectoral (source0.tex lines 211--219). The article's complete original proof of it is delivered with it (source0.tex lines 221--250, one \texttt{proof} environment of 30 lines). No mathematics is written, repaired or re-derived: the statement and the proof are the article's own lines, in the article's own notation, and no retained line is substituted or re-flowed.  Retained uncounted as the source-local context the proof needs: lines 84--87.  Nothing was omitted from any retained span, so the package carries no omission record.  No retained line contains a citation, so the wrapper carries no bibliography and no bibliography entry.  No retained line contains a figure environment, an image-inclusion command or drawing code, so no image is delivered and the package declares no asset dependency.  Editorial formatting only, in addition: an AMS \texttt{amsart} preamble that restates the article's own declarations and macros used by the retained lines, namely the environments \texttt{theorem} on separate counters, together with the standard packages those lines need.  The retained environments print in this document as Theorem 1 in order of appearance, whereas the article prints the same items with the numbering of its own full document.  The complete native source of the article as distributed in the arXiv e-print for this exact version is bundled unchanged as \texttt{sources/191397/source0.tex}.
\begin{equation}
\frac{1}{\sin\theta}\frac{d}{d\theta}\left(\sin\theta\frac{d\Theta}{d\theta}\right) + \left[\nu(\nu+1) - \frac{m^2}{\sin^2\theta}\right]\Theta = 0,
\label{eq:angular_ode}
\end{equation}

\begin{theorem}[Sectoral Solution]
\label{thm:sectoral}
For $\nu = m$ with $m$ any positive real number, the function
\begin{equation}
\Theta(\theta) = (\sin\theta)^m
\label{eq:sectoral_solution}
\end{equation}
is an exact solution of the associated Legendre equation~\eqref{eq:angular_ode}. This solution vanishes at both poles $\theta = 0$ and $\theta = \pi$ and is everywhere positive on the open interval $(0, \pi)$.
\end{theorem}

\begin{proof}
The proof proceeds by direct substitution. Writing $\Theta = \sin^m\theta$ for notational brevity, we compute the required derivatives:
\begin{align}
\frac{d\Theta}{d\theta} &= m\sin^{m-1}\theta\cos\theta, \\
\sin\theta\frac{d\Theta}{d\theta} &= m\sin^m\theta\cos\theta, \\
\frac{d}{d\theta}\left(\sin\theta\frac{d\Theta}{d\theta}\right) &= m\frac{d}{d\theta}(\sin^m\theta\cos\theta) \notag \\
&= m\left[m\sin^{m-1}\theta\cos^2\theta - \sin^{m+1}\theta\right].
\end{align}
The first term in equation~\eqref{eq:angular_ode} therefore becomes
\begin{equation}
\frac{1}{\sin\theta}\frac{d}{d\theta}\left(\sin\theta\frac{d\Theta}{d\theta}\right) = m^2\sin^{m-2}\theta\cos^2\theta - m\sin^m\theta.
\label{eq:first_term}
\end{equation}

For the potential term, we set $\nu(\nu+1) = m(m+1)$ and compute
\begin{align}
\left[\nu(\nu+1) - \frac{m^2}{\sin^2\theta}\right]\Theta &= \left[m(m+1) - \frac{m^2}{\sin^2\theta}\right]\sin^m\theta \notag \\
&= m(m+1)\sin^m\theta - m^2\sin^{m-2}\theta.
\label{eq:potential_term}
\end{align}

Adding equations~\eqref{eq:first_term} and~\eqref{eq:potential_term}:
\begin{align}
&m^2\sin^{m-2}\theta\cos^2\theta - m\sin^m\theta + m(m+1)\sin^m\theta - m^2\sin^{m-2}\theta \notag \\
&= m^2\sin^{m-2}\theta(\cos^2\theta - 1) + m^2\sin^m\theta \notag \\
&= -m^2\sin^{m-2}\theta \cdot \sin^2\theta + m^2\sin^m\theta \notag \\
&= -m^2\sin^m\theta + m^2\sin^m\theta = 0.
\end{align}
The function $(\sin\theta)^m$ thus satisfies equation~\eqref{eq:angular_ode} identically when $\nu = m$, completing the proof.
\end{proof}

\end{document}
